% !TeX root = ../main.tex
% Add the above to each chapter to make compiling the PDF easier in some editors.

\chapter{Future Work}\label{chap:future_work}

The practical analysis highlights a broader limitation in the current tensor contraction ecosystem. Many state-of-the-art systems achieve strong performance by tightly coupling contraction path optimization, layout decisions, contraction execution, and low-level compilation. While this produces highly optimized implementations, it also makes it difficult for new algorithms to reuse only the lower-level performance infrastructure. As a result, researchers proposing new path-selection or permutation-optimization algorithms may need to reimplement large parts of the execution stack before their ideas can be evaluated fairly in practice, as with our attempt here.

A promising direction for future work is the development of a tensor network compiler. Such a compiler would separate high-level contraction optimization from low-level code generation and memory management. In this model, algorithms such as the proposed greedy permutation strategy would output an intermediate representation describing the contraction path, tensor layouts, frozen tensors, permutation constraints, and memory-reuse opportunities. The compiler would then lower this representation into optimized kernels, memory plans, and execution schedules. This would allow new algorithms to benefit from the same low-level optimizations currently available in mature libraries without requiring each algorithm to reimplement the complete stack.

This separation would also make practical evaluation more reliable. If different algorithms target the same compiler backend, their differences can be measured at the algorithmic level rather than being dominated by differences in implementation quality. Additionally, a compiler-based approach could introduce advanced features such as allocator-aware contraction scheduling, automatic view reuse, memory-pool planning, fusion of permutation and contraction operations, and hardware-specific kernel generation. These optimizations are difficult to implement independently inside every new algorithm, but they are natural responsibilities for a shared compiler infrastructure.

In addition to structural ecosystem improvements, several direct enhancements to the core algorithm are required to fully realize its potential. First, the greedy permutation strategy must be extended to incorporate more sophisticated leaf node handling, ensuring that initial tensor layouts and their immediate transformations are optimized globally rather than myopically. Second, and perhaps more critically, the algorithm needs to be made natively memory-aware. Currently, the strategy operates under a memory-agnostic paradigm, occasionally suggesting permutation schedules that can exceed strict hardware memory constraints during execution. Managing these limits is presently deflected to the user, who must manually determine how many tensors to freeze or exclude from permutation. Future iterations of the algorithm should natively incorporate these memory constraints into the optimization objective, automatically introducing strategic slicing and tensor freezing whenever a path risks breaching the available memory budget. Integrating these layout and capacity constraints directly into the search heuristic will remove manual tuning and guarantee out-of-the-box feasibility for large-scale tensor networks.

Therefore, the main future direction suggested by this work is not only to improve the greedy strategy itself, but to improve the surrounding ecosystem in which such strategies are executed. A modular tensor network compiler would make the field more extensible, more comparable, and more practically useful. It would allow researchers to focus on new contraction and layout ideas while still benefiting from state-of-the-art low-level optimization. In turn, this could help close the gap between theoretical memory savings and practical end-to-end performance.